\section{Experiments}
\label{sec:experiments}

\subsection{Test Environment}

All benchmarks were executed on a single host machine with the specifications listed in
Table~\ref{tab:hardware}. The three database containers and the benchmark client
shared this machine with no concurrent workloads running during test sessions.

\begin{table*}[t]
  \centering
  \caption{Test environment hardware and software versions.}
  \label{tab:hardware}
  \begin{tabular}{ll}
    \hline
    \textbf{Component} & \textbf{Details} \\
    \hline
    CPU       & Intel Core i5-11400F (6c/12t, 2.60--4.40\,GHz) \\
    RAM       & 32\,GB DDR4 \\
    Disk      & 500\,GB NVMe SSD \\
    OS        & NixOS (Linux 6.12) \\
    \hline
    IoTDB       & \texttt{apache/iotdb:1.3.0-standalone} \\
    InfluxDB    & \texttt{influxdb:2.7-alpine} \\
    TimescaleDB & \texttt{timescale/timescaledb:latest-pg15} \\
    Benchmark   & iot-benchmark~\cite{iotbenchmark} (Java, Maven) \\
    \hline
  \end{tabular}
\end{table*}

CPU utilization figures are expressed as a percentage of a single logical core, consistent
with \texttt{docker stats} reporting conventions. A value of 600\% would indicate full
utilization of all six physical cores; single-database values in the 20--100\% range
indicate load compatible with the available compute budget.

\subsection{Write Performance}
\label{sec:write-results}

\subsubsection{Baseline Sequential Ingestion}

Table~\ref{tab:write-baseline} shows baseline write performance at small scale under the tuned
configuration. IoTDB achieves 2.16\,M\,pts/s, approximately 5$\times$ InfluxDB and 12$\times$
TimescaleDB.

\begin{table*}[t]
  \centering
  \caption{Baseline write performance at small scale (tuned). Throughput in pts/s.}
  \label{tab:write-baseline}
  \begin{tabular}{lrrrr}
    \hline
    \textbf{DB} & \textbf{pts/s} & \textbf{Avg\,Lat} & \textbf{P99\,Lat} & \textbf{Avg\,CPU} \\
    \hline
    IoTDB       & 2{,}163{,}809 &  1.73\,ms &   1.18\,ms &  5.1\% \\
    InfluxDB    &   397{,}894   & 12.00\,ms & 190.50\,ms &  8.5\% \\
    TimescaleDB &   174{,}787   & 27.70\,ms & 223.00\,ms & 22.5\% \\
    \hline
  \end{tabular}
\end{table*}

The results are consistent with three compounding architectural factors.

\textbf{Protocol depth.} iot-benchmark connects to IoTDB via the Thrift RPC Session API in SESSION\_BY\_TABLET mode,
submitting typed columnar tablets directly over a persistent TCP connection. One
syscall submits a full batch directly into the TsFile write buffer. InfluxDB requires an
HTTP request per write, including connection management, line-protocol text parsing, and
TSI index lookup. TimescaleDB routes every write through the full PostgreSQL stack:
TCP, SQL parsing, query planning, hypertable routing, WAL write, and MVCC tuple formatting.
Each layer adds latency; the PostgreSQL stack adds the most.

\textbf{Memory buffering and durability trade-offs.} IoTDB accumulates tablets in its
JVM heap and acknowledges writes before disk persistence, flushing to TsFile asynchronously
when the memtable fills. This explicitly trades durability guarantee for throughput, a
trade-off that must be evaluated per deployment. InfluxDB's TSM cache operates similarly
but is bounded (8\,GB in the tuned configuration). TimescaleDB follows standard PostgreSQL
semantics: every write is flushed through the WAL before acknowledgement, providing the
strongest durability guarantee at the cost of per-write synchronous input/output.

\textbf{Index maintenance.} IoTDB does not maintain synchronous indexes on the write path;
data goes directly to the in-memory buffer without updating auxiliary structures.
InfluxDB updates the Time Series Index (TSI) on every write. TimescaleDB maintains
B-tree and BRIN indexes that grow with data volume, and the PostgreSQL autovacuum process,
responsible for reclaiming space from obsolete records, adds overhead
as the hypertable expands.

The combination of these three factors suggests that IoTDB's write advantage has an
architectural basis: the binary protocol, asynchronous buffering, and absence of
synchronous index maintenance would act cumulatively and are not eliminated by configuration
tuning.

\begin{figure}[ht]
  \centering
  \includegraphics[width=\columnwidth]{figures/en_us/chart_write_throughput_small}
  \caption{Write throughput (pts/s, log scale) for all four write workloads at small
  scale (tuned). IoTDB leads across every test type.}
  \label{fig:write_throughput_small}
\end{figure}

\subsubsection{Batch Size Sensitivity}

Table~\ref{tab:batch} shows how throughput changes between BATCH\_SIZE=1 (one write request
per point) and BATCH\_SIZE=1000 (1,000 points per request).

\begin{table}[ht]
  \centering
  \caption{Batch size sensitivity at small scale (tuned). Throughput in pts/s.}
  \label{tab:batch}
  \begin{tabular}{lrrr}
    \hline
    \textbf{DB} & \textbf{BATCH=1} & \textbf{BATCH=1000} & \textbf{Ratio} \\
    \hline
    IoTDB       &  28{,}593 & 11{,}827{,}572 & 413$\times$ \\
    InfluxDB    &   5{,}024 &  2{,}667{,}256 & 531$\times$ \\
    TimescaleDB &   5{,}120 &    214{,}927   &  42$\times$ \\
    \hline
  \end{tabular}
\end{table}

IoTDB and InfluxDB show 413--531$\times$ speedup, reflecting the high cost of per-request
overhead. For InfluxDB, this is an HTTP round-trip plus TCP connection overhead; for IoTDB,
it is the fixed cost of a Thrift call per tablet. TimescaleDB's 42$\times$ ratio is
markedly lower because PostgreSQL's WAL flush already dominates at BATCH=1: the per-row
overhead is so high that larger batches provide proportionally less benefit; the bottleneck
shifts from protocol overhead to storage input/output.

This result has direct practical significance. An IoT application that sends one sensor
reading per HTTP request to InfluxDB will achieve 531$\times$ less throughput than one
that buffers locally and writes in batches. For InfluxDB and IoTDB, aggregating writes into
batches is the highest-impact optimization available.

\subsubsection{Out-of-Order Ingestion}

Table~\ref{tab:ooo} shows how each database handles the 20\% late-arriving data workload.

\begin{table*}[t]
  \centering
  \caption{Out-of-order penalty: throughput (pts/s) and relative degradation
  from baseline write, small scale (tuned).}
  \label{tab:ooo}
  \begin{tabular}{lrrr}
    \hline
    \textbf{DB} & \textbf{WRITE\,(pts/s)} & \textbf{OOO\,(pts/s)} & \textbf{Penalty} \\
    \hline
    IoTDB       & 2{,}163{,}809 & 1{,}871{,}351 & $-$13.5\% \\
    InfluxDB    &   397{,}894   &   255{,}634   & $-$35.7\% \\
    TimescaleDB &   174{,}787   &   160{,}121   & $-$8.4\% \\
    \hline
  \end{tabular}
\end{table*}

TimescaleDB shows the smallest relative degradation, a result consistent with the
order-independence of PostgreSQL's B-tree index updates. IoTDB's 13.5\% penalty
is consistent with the cost of compacting disordered entries in the LSM-tree to restore
the sorted order required by TsFile chunks. InfluxDB's largest relative penalty (35.7\%) is
consistent with TSM block rewrites needed to accommodate out-of-sequence points.

In absolute throughput, IoTDB still far leads under out-of-order conditions (1.87\,M vs.\
255\,K pts/s for InfluxDB). Two dimensions must therefore be considered separately: absolute
out-of-order throughput (IoTDB leads at every scale) and relative degradation from write
disorder (TimescaleDB is most resilient with the smallest fractional penalty).

\begin{figure}[ht]
  \centering
  \includegraphics[width=\columnwidth]{figures/en_us/chart_ooo_degradation}
  \caption{Out-of-order throughput degradation (\%) relative to baseline sequential write,
    across all three scales. TimescaleDB's B-tree remains stable; InfluxDB's TSM
  block rewrites consistently incur the largest relative penalty.}
  \label{fig:ooo_degradation}
\end{figure}

\subsection{Read Performance}
\label{sec:read-results}

Read throughput in pts/s varies with the number of points returned per query, making
that value incomparable across different read types. Latency is used as the primary
metric in this section. It is also worth noting that read tests run sequentially
on the same dataset, so later tests benefit from warmer OS page caches. When that
cache warm-up is itself the mechanism behind a database's performance, as with
TimescaleDB's \texttt{shared\_buffers} in downsampling, the result is expected and
reproducible in production (see Section~\ref{sec:limitations}).

\subsubsection{Latest-Point Queries}

Table~\ref{tab:latest} shows latest-point average latency across all three scales for the
tuned runs.

\begin{table}[ht]
  \centering
  \caption{Latest-point average latency (ms) across scales (tuned). Lower is better.}
  \label{tab:latest}
  \begin{tabular}{lrrr}
    \hline
    \textbf{DB} & \textbf{Small} & \textbf{Medium} & \textbf{Large} \\
    \hline
    TimescaleDB &  0.38 &  0.41 &  0.62 \\
    IoTDB       &  1.24 &  0.51 &  0.68 \\
    InfluxDB    &  9.24 & 21.72 & 23.24 \\
    \hline
  \end{tabular}
\end{table}

TimescaleDB and IoTDB maintain sub-millisecond or near-sub-millisecond latency across
a hundredfold increase in data volume. Both exploit database-specific shortcuts for last-value retrieval.
TimescaleDB's behavior is consistent with its time-ordered hypertable partitioning~\citep{timescaledb},
where latest-point queries need only access the most recent chunk,
regardless of total stored volume. IoTDB maintains a dedicated in-memory last-value cache per series, serving
the query entirely from RAM without accessing TsFile~\citep{iotdblastquery}.

InfluxDB has no equivalent shortcut. It must seek to the tail of each series' TSM file
blocks, and that seek cost grows as more TSM files accumulate over the dataset lifetime.
The resulting progression (9.24\,ms at small, 21.72\,ms at medium, 23.24\,ms at
large) confirms an approximately linear relationship between stored data and latest-point
cost. Across the evaluated scales, InfluxDB's latest-point latency showed approximately linear
growth with stored data volume. In deployments without active retention or downsampling
policies, this pattern may make it disadvantageous for high-frequency real-time
monitoring; deployments that apply such policies may mitigate this behavior.

\begin{figure}[ht]
  \centering
  \includegraphics[width=\columnwidth]{figures/en_us/chart_latest_point_scaling}
  \caption{Latest-point average latency (ms) across scales. TimescaleDB and IoTDB stay
    near zero (O(1) access paths); InfluxDB grows linearly as the TSM file count
  increases.}
  \label{fig:latest_point_scaling}
\end{figure}

\subsubsection{Downsampling}

Table~\ref{tab:downsample} shows GROUP~BY time-bucket downsampling latency.

\begin{table}[ht]
  \centering
  \caption{Downsampling average latency (ms) across scales (tuned). Lower is better.}
  \label{tab:downsample}
  \begin{tabular}{lrrr}
    \hline
    \textbf{DB} & \textbf{Small} & \textbf{Medium} & \textbf{Large} \\
    \hline
    TimescaleDB & 0.61 & 0.52 & 0.73 \\
    IoTDB       & 1.81 & 0.64 & 0.95 \\
    InfluxDB    & 5.71 & 5.65 & 9.22 \\
    \hline
  \end{tabular}
\end{table}

TimescaleDB's \texttt{time\_bucket()} aggregation stays below 1\,ms across a hundredfold
increase in data volume, a result that holds consistently across all tested scales. This result is consistent with TimescaleDB's hypertable chunk exclusion mechanism~\citep{timescaledb}, by which only chunks that overlap the queried
time window are accessed, regardless of total stored volume. The warmed \texttt{shared\_buffers} in
the tuned runs reinforces this behavior, keeping relevant chunks in memory.
This combination makes TimescaleDB suitable for dashboard workloads that
consistently query recent time windows.

IoTDB's downsampling latency improves from small to medium scale (1.81\,ms $\to$ 0.64\,ms)
as the longer medium-scale run gives the JVM JIT compiler time to optimise the chunk
aggregation hot path~\citep{barrett2017warmup}. At large scale it remains under 1\,ms. InfluxDB incurs higher per-query overhead (5--9\,ms).

\subsubsection{Value-Filter Queries}

Value-filter queries (predicates on sensor values such as \texttt{temperature > 85.0})
are the most architecturally discriminating workload in this study. Time-based indexes
cannot answer them directly; each engine must either scan all stored values or exploit
a secondary data structure to prune non-matching storage blocks.

Table~\ref{tab:vfilter} shows value-filter latency and its scaling behavior across scales.

\begin{table*}[t]
  \centering
  \caption{Value-filter average latency (ms) across scales (tuned).
  The ratio column shows latency growth from small to large scale (data volume one hundred times larger).}
  \label{tab:vfilter}
  \begin{tabular}{lrrrr}
    \hline
    \textbf{DB} & \textbf{Small} & \textbf{Medium} & \textbf{Large} & \textbf{S$\to$L ratio} \\
    \hline
    IoTDB       &   3.72 &    9.80 &   19.42 &  5.2$\times$ \\
    TimescaleDB &  23.98 &  130.08 &  510.05 & 21.3$\times$ \\
    InfluxDB    &  71.09 &  310.75 &  496.81 &  7.0$\times$ \\
    \hline
  \end{tabular}
\end{table*}

IoTDB's advantage is large and improves with scale. Each TsFile chunk stores
per-series min/max statistics. A value predicate can skip entire chunks whose stored range
does not intersect the predicate threshold, providing lightweight secondary indexing at
the storage layer. At large scale, with more data per chunk, a larger fraction of chunks
can be pruned, producing the sublinear scaling ratio (5.2$\times$ for a hundredfold increase in data).
This is the expected behavior of a statistics-based skip structure: the pruning ratio
improves as chunks grow denser.

TimescaleDB uses PostgreSQL's BRIN (Block Range INdex), which records min/max values per
block range on disk. BRIN provides coarser-grained pruning than IoTDB's per-chunk statistics
and degrades at a superlinear rate (21.3$\times$ for a hundredfold increase in data). At large scale,
TimescaleDB averages 510\,ms per value-filter query, a 26$\times$ disadvantage relative
to IoTDB. This workload also drives TimescaleDB to its highest CPU utilisation (77\% average,
83\% peak), indicating the query executor is the bottleneck.

InfluxDB has no statistics-based pruning structure but benefits from some coarse file-level
skipping based on time-range metadata in the TSM index, producing a 7$\times$ scaling ratio,
significantly better than TimescaleDB's but still 10$\times$ worse than IoTDB in absolute
terms at large scale.

\begin{figure}[ht]
  \centering
  \includegraphics[width=\columnwidth]{figures/en_us/chart_value_filter_scaling}
  \caption{Value-filter average latency (ms, log scale) across scales. IoTDB's chunk
    min/max statistics produce sublinear growth (5.2$\times$ for a hundredfold increase in data);
  TimescaleDB's BRIN index degrades superlinearly (21.3$\times$).}
  \label{fig:value_filter_scaling}
\end{figure}

\subsubsection{Range Queries}

Range queries retrieve all data points within a time window without value filtering.
TimescaleDB is fastest here: 0.54\,ms at small scale and 0.66\,ms at large, benefiting
from hypertable chunk pruning (only chunks overlapping the query window are opened) and the
warm OS page cache that avoids disk input/output for frequently queried recent time windows. IoTDB
achieves 1.19--1.21\,ms and is effectively scale-insensitive within this data range.
InfluxDB achieves 5.41--8.68\,ms, where the TSM block scanner bears higher per-file
overhead than TimescaleDB's chunk-level sequential scan.

\subsection{Memory Footprint}
\label{sec:memory-results}

Memory measurements require careful interpretation because the original shared-start
benchmark runs produced substantially inflated figures through a methodological artefact.

\subsubsection{TimescaleDB}

TimescaleDB reported approximately 5.4\,GB from the very first test in the
original runs. This was not due to PostgreSQL's buffer pool, but rather to the operating system's page cache, which had been accumulated during IoTDB's preceding write sessions. Linux accounts for file-system page cache
toward a container's RSS in \texttt{docker stats}, and the idle TimescaleDB container
had been running for hours while IoTDB generated gigabytes of input/output. In the tuned
sequential runs, TimescaleDB starts at 239\,MB and grows to approximately 1.8\,GB
after the write tests warm \texttt{shared\_buffers}. The 1.8\,GB figure is the
accurate operational footprint for this dataset size under the tuned configuration.

\subsubsection{IoTDB}

IoTDB reported approximately 12.5\,GB flat throughout the original runs.
During the hours it spent idle while InfluxDB and TimescaleDB were tested, the JVM
expanded the heap up near the configured maximum and did not return those pages to the OS,
which is standard behavior of the JVM heap manager. In the tuned runs, the JVM starts
at 1.7\,GB and grows to approximately
10\,GB by the end of the read phase, reflecting actual working-set demand per workload
phase.

\subsubsection{InfluxDB}

InfluxDB is the most memory-efficient across write workloads: 190--940\,MB in
tuned write tests at small scale. However, this advantage does not hold unconditionally.
At large scale, BATCH-LARGE peaks at 2.9\,GB as the TSM write cache holds large volumes
of pending batches alongside concurrent compaction buffers. The default 1\,GB limit in InfluxDB applies only to the TSM cache ceiling. Total process RSS includes the series index, WAL buffers, and Go runtime overhead, components that together can exceed that limit under high write load.

The practical implication for deployment planning is that IoTDB's dominant write throughput
requires a substantial memory commitment ($\sim$10\,GB at medium scale, growing to $\sim$15\,GB
at large). For memory-constrained environments (IoT edge gateways, embedded systems,
or cloud instances with limited RAM), InfluxDB's leaner footprint may be the deciding
factor even where write throughput is the primary concern.

\begin{figure}[ht]
  \centering
  \includegraphics[width=\columnwidth]{figures/en_us/chart_memory_scaling}
  \caption{Peak memory footprint (MB) across scales (tuned runs). IoTDB's JVM working
    set grows to 15\,GB at large scale. InfluxDB remains the most
  memory-efficient option at every scale.}
  \label{fig:memory_scaling}
\end{figure}

\subsection{Scale Effects}
\label{sec:scale-effects}

\begin{figure*}[t]
  \centering
  \includegraphics[width=\textwidth]{figures/en_us/chart_heatmap_latency}
  \caption{Latency degradation heatmap across read workload types and scale transitions
    (tuned runs). Green indicates stable latency; red indicates high growth.
    IoTDB is consistently green except for value-filter (small$\to$medium);
  TimescaleDB shows severe degradation on value-filter at both transitions.}
  \label{fig:heatmap_latency}
\end{figure*}

\subsubsection{Write Throughput Scaling}

Table~\ref{tab:write-scale} shows how baseline write throughput evolves across scales
in the tuned runs.

\begin{table}[ht]
  \centering
  \caption{Baseline write throughput (pts/s) across scales (tuned).}
  \label{tab:write-scale}
  \begin{tabular}{lrrr}
    \hline
    \textbf{DB} & \textbf{Small} & \textbf{Medium} & \textbf{Large} \\
    \hline
    IoTDB       & 2{,}163{,}809 &  7{,}991{,}224 & 23{,}607{,}475 \\
    InfluxDB    &   397{,}894   &    392{,}975   &  1{,}256{,}602 \\
    TimescaleDB &   174{,}787   &    175{,}685   &    400{,}455   \\
    \hline
  \end{tabular}
\end{table}

IoTDB's write throughput scales superlinearly: 2.2\,M pts/s at small, 8.0\,M at medium,
23.6\,M at large. This growth is consistent with the JVM JIT compilation warm-up effect:
short benchmarks do not give the
JIT sufficient time to fully optimise all hot paths~\citep{barrett2017warmup}. At medium and large scale the longer
runs would allow the JIT to reach peak compilation, saturating the Thrift write path. Small-scale
IoTDB write numbers therefore likely understate its steady-state ceiling; the large-scale figure
is the most representative of sustained production throughput.

InfluxDB remains roughly constant between small and medium, a pattern suggestive of
saturation by WAL and TSM compaction regardless of volume. At large
scale, the observed gain is consistent with amortization of compaction overhead over larger volumes.
TimescaleDB shows a similar pattern: flat from small to medium,
then improving at large scale, suggestive of better use of checkpoint intervals
with longer runs.

\begin{figure}[ht]
  \centering
  \includegraphics[width=\columnwidth]{figures/en_us/chart_write_scaling}
  \caption{Baseline write throughput (pts/s, log scale) across scales (tuned). IoTDB's
    superlinear growth reflects JVM JIT saturation; InfluxDB and TimescaleDB
  remain flat until large scale.}
  \label{fig:write_scaling}
\end{figure}

\subsubsection{IoTDB BATCH-SMALL JIT Saturation}

IoTDB's \texttt{batch-small} throughput across scales (28\,K $\to$ 81\,K $\to$ 446\,K pts/s)
is one of the clearest demonstrations of JIT-driven scaling in the dataset. Average latency
simultaneously decreases from 1.22\,ms to 0.40\,ms as scale increases, which is
only possible if the improvement comes from code optimization rather than reduced load.
The large-scale figure (446\,K pts/s, 0.40\,ms average latency) represents IoTDB's
steady-state single-point ingestion throughput once the JIT has fully compiled the
Thrift serialisation hot path~\citep{barrett2017warmup}.

\subsubsection{TimescaleDB BATCH-LARGE Timeout}

TimescaleDB's \texttt{batch-large} test hits the 3{,}600-second benchmark timeout at large
scale in both the original and tuned runs. With 50 devices and a batch size of 1{,}000 rows,
each write request touches 50 hypertable partitions simultaneously and triggers up to 50
concurrent WAL flushes. The cumulative input/output cost exhausts the time budget before the full
loop count completes. This does not represent an anomaly; it is the expected outcome of partition fan-out in the hypertable under high-concurrency, high-volume write scenarios. For bulk-loading
workloads targeting TimescaleDB at scale, the partition count must be weighed against the
WAL flush amplification.

\subsubsection{InfluxDB Latest-Point Linear Growth}

Among the multi-scale trends, InfluxDB's latest-point latency progression (9.24 $\to$
21.72 $\to$ 23.24\,ms) deserves particular attention because it reveals a structural
scalability issue. Every incremental TSM file written to disk adds to the seek depth
that a latest-point query must traverse. TimescaleDB and IoTDB show flat or even
decreasing latest-point latency (due to IoTDB's in-memory cache~\citep{iotdblastquery} and JIT warm-up~\citep{barrett2017warmup})
across the same hundredfold increase in data volume. Any deployment that will accumulate months or years of
data will see InfluxDB's latest-point cost continue to grow, while TimescaleDB's
and IoTDB's costs remain optimized.

\subsection{Impact of Configuration Tuning}
\label{sec:tuning-results}

\subsubsection{TimescaleDB Gains from pgtune}

Table~\ref{tab:tsdb-tuning} compares TimescaleDB write performance before and after
pgtune at small scale.

\begin{table*}[t]
  \centering
  \caption{TimescaleDB write performance: original vs.\ tuned (small scale).
  $\Delta$ shows percentage change; negative means improvement for latency.}
  \label{tab:tsdb-tuning}
  \begin{tabular}{lrrrr}
    \hline
    \textbf{Test} & \textbf{Old pts/s} & \textbf{New pts/s} & $\Delta$tput & $\Delta$lat \\
    \hline
    WRITE        & 133{,}677 & 174{,}787 & $+$30\% & $-$25\% \\
    BATCH-LARGE  & 160{,}546 & 214{,}927 & $+$34\% & $-$25\% \\
    OUT-OF-ORDER & 132{,}189 & 160{,}121 & $+$21\% & $-$21\% \\
    \hline
  \end{tabular}
\end{table*}

The primary driver is \texttt{shared\_buffers=8GB}, which provides PostgreSQL a large
enough buffer pool to avoid buffer replacement that would force extra WAL writes.
\texttt{work\_mem=27413kB} eliminates disk writes during sort-based filtering, which explains the 40\% reduction in value-filter latency -- a workload that requires sorting filtered results.
\texttt{checkpoint\_completion\_target=0.9} spreads checkpoint input/output more
evenly across the checkpoint interval, reducing the write stall spikes that inflate P99
latency (BATCH-LARGE P99 dropped from 1{,}069\,ms to 812\,ms). These improvements
indicate that default PostgreSQL settings are intentionally conservative to ensure
correctness on low-resource hardware, and that tuning via pgtune or an equivalent tool
should be applied before reporting results on mid- to high-end hardware.

\subsubsection{InfluxDB Gains from Isolation}

InfluxDB received no configuration changes in the tuned runs. Its improvements come entirely
from running without competing databases. All read latencies improved 15--21\%, consistent
with better TSM file cache utilization once TimescaleDB's page cache no longer competes for
OS buffer space. The most dramatic result was BATCH-LARGE P99: 205\,ms $\to$ 60\,ms
($-$71\%). In the original setup, TimescaleDB was writing to disk during
InfluxDB's test session, periodically starving InfluxDB's TSM compaction engine of input/output
bandwidth. These rare but severe input/output stalls produced extreme latency spikes that, while not
affecting average throughput, considerably inflated P99. Isolation removed these stalls.

\subsubsection{Correcting Two Large-Scale False Positives}

The original large-scale runs produced two findings that the tuned reruns reversed.

\textbf{TimescaleDB out-of-order penalty.} The original large-scale run showed a 23\%
OOO throughput penalty for TimescaleDB, attributed to decompression-recompression cycles
from out-of-order inserts into older compressed chunks. The tuned rerun shows only 2.7\%
degradation, consistent with the pattern at small and medium scale. The most likely hypothesis: out-of-order inserts into older chunks require reading
from disk for decompression. In the original shared-start session, IoTDB's expanding JVM heap (above
15\,GB) likely displaced the OS page cache pages holding those chunks, forcing cold reads
on every OOO insert into an older chunk. With exclusive access, the chunks would
remain warm in cache and decompression costs would be negligible.

\textbf{TimescaleDB CPU saturation.} The original READ workload at large scale averaged
97.7\% CPU utilization with a P99 latency of 4{,}232\,ms, leading to the conclusion that
TimescaleDB ``runs out of compute'' and could not be improved without a stronger CPU or
horizontal sharding. The tuned rerun shows 77.4\% average CPU and 1{,}850\,ms P99.
The original 97.7\% figure resulted from all three databases and the benchmark client simultaneously competing for the same six CPU cores. TimescaleDB's performance was therefore not limited by the hardware's physical capacity, but by excessive processing contention caused by concurrent processes.

These two reversals underscore the importance of isolated, correctly-tuned benchmark
methodology. Without sequential container lifecycles and standard database configuration,
both claims would have been published as validated findings.

\begin{figure*}[t]
  \centering
  \includegraphics[width=\textwidth]{figures/en_us/chart_read_saturation_comparison}
  \caption{READ workload average CPU utilisation (top) and P99 latency (bottom, log
    scale) for original (hatched) vs.\ tuned (solid) runs across all three
    scales. TimescaleDB's large-scale CPU drops from 98\% to 77\% and P99 from
  4{,}233\,ms to 1{,}850\,ms once competing containers are removed.}
  \label{fig:read_saturation_comparison}
\end{figure*}

\subsection{Mixed Read Workload}

The \texttt{read} workload runs all ten iot-benchmark query types with equal weight,
providing a composite read baseline that includes range queries, aggregations, downsampling,
latest-point lookups, and value-predicate queries simultaneously. At small scale (tuned),
TimescaleDB averages 0.67\,ms, IoTDB averages 2.11\,ms, and InfluxDB averages 14.28\,ms
per operation across the combined set. These figures are consistent with the individual
query results presented above: TimescaleDB dominates for time-windowed queries, which
constitute the majority of the ten types, while IoTDB's advantage on value-filter queries
is diluted by the equal weighting.

The mixed workload's primary analytical role in this study is as a tuning benchmark:
TimescaleDB's read latency improved 38\% after pgtune, confirming that multiple
query types benefit from a correctly sized buffer pool, not only the value-filter.

\subsection{Summary}
\label{sec:results-summary}

Table~\ref{tab:summary} consolidates the workload-winner mapping across all tested scenarios.

\begin{table*}[t]
  \centering
  \caption{Workload recommendation summary across all tested scenarios.}
  \footnotesize
  \label{tab:summary}
  \begin{tabular}{l l p{0.55\linewidth}}
    \hline
    \textbf{Workload} & \textbf{Best DB} & \textbf{Key architectural reason} \\
    \hline
    Sequential ingestion (ordered)     & IoTDB       & Binary tablet protocol + in-memory buffering before TsFile flush \\
    Bulk loading (BATCH=1000)          & IoTDB       & Columnar vectorised write path; 11.8\,M pts/s at small scale \\
    Single-point ingestion (BATCH=1)   & IoTDB       & Lower per-call Thrift overhead vs HTTP round-trip \\
    Out-of-order (absolute throughput) & IoTDB       & 1.87\,M pts/s OOO vs 255\,K (InfluxDB); leads at every scale \\
    Out-of-order (relative tolerance)  & TimescaleDB & B-tree updates are order-independent \\
    Dashboard downsampling             & TimescaleDB & \texttt{time\_bucket()} in hot \texttt{shared\_buffers}; constant latency \\
    Bulk time-range export             & TimescaleDB & Chunk pruning + warm OS page cache \\
    Threshold alerting (value-filter)  & IoTDB       & Chunk min/max statistics; sublinear scaling with data volume \\
    Real-time last-value lookup        & TimescaleDB & Sub-ms reverse B-tree scan on most recent chunk \\
    Memory efficiency (write)          & InfluxDB    & Bounded TSM cache; no JVM pre-allocation \\
    \hline
  \end{tabular}
\end{table*}

IoTDB is the clear choice for write-dominated workloads but requires substantial memory
($\sim$10\,GB at medium scale, growing to $\sim$15\,GB at large) and does not support
relational joins. TimescaleDB trades raw write speed for SQL compatibility, the full
PostgreSQL ecosystem, and consistently strong read performance for time-windowed queries
once data is cached. InfluxDB occupies the middle ground: better write throughput than
TimescaleDB, better memory efficiency than IoTDB, but weaker on complex value-predicate
queries and latest-value lookups at scale. No single database dominates across all workloads.
